\chapter{Mental models, reasoning and rationality}
\label{sec:p19-rationality-2-3}

A mental-model theory is a theory of how people represent situations and
draw conclusions from them, not a replacement name for model-theoretic
semantics. Johnson-Laird's 2021 presentation treats deduction, assessments
of probability, induction and explanation as related uses of models. Its
central proposal is that a reasoner constructs representations of what
could happen, draws a conclusion and, when necessary, searches for an
alternative representation that would defeat it. Limited representation
and search can produce predictable mistakes. Accurate models, informative
conclusions and explanatory power are the proposed aims of right reasoning
\cite[pp.~147--156]{JohnsonLaird2021MentalModels}.

Three questions must remain separate. Does a conclusion satisfy a specified
formal standard? What information did a participant interpret the premises
as providing? What cognitive process produced the response? A truth table
can answer the first question under a stipulated interpretation without
answering the other two. Conversely, a response made by most participants
does not establish a valid inference or a sound normative requirement.
\Cref{def:p19-empirical-question,def:p19-evidence-record} gives the evidence
framework used here. The conditional and conjunction treatments at
\cref{sec:p19-conditional-reasoning-suppression,sec:p19-conjunction-judgments}
provide related task-specific contrasts.

\section{Representations, standards and the mental-logic comparison}
\label{sec:p19-r23-representation}

The source calls a model \emph{iconic} when relations in the representation
correspond to relations in what it represents. A spatial arrangement can
represent who is beside whom, and a changing arrangement can represent a
sequence of movements. Iconicity is not a requirement that every model
be a conscious picture. A sentential possibility can be represented by
tokens, and a magnitude of belief by an analog scale. The theory also
postulates symbolic operations, including negation. Its intuitive process
represents what is true in a possibility while often omitting what is
false; deliberation can make the omitted alternatives explicit. This is a
psychological proposal about representations and their use, not a theorem
that a particular brain implements a valuation list.

The source begins with a 2021 political illustration: a disjunctive threat
of tariffs against Russia, China or both licenses, in an ordinary reading,
the thought that tariffs against Russia are possible. Its function is to
illustrate an interpretation of a disjunction, not to assert a current
political event. The same issue is made precise below with neutral atoms.
The source's evolutionary account suggests that an ability to represent
the environment and simulate its changes is older than explicit formal
calculi. Such a proposed genealogy does not establish either the truth
of every representation or the adequacy of one particular cognitive model.

\emph{Mental logic}, in this comparison, refers to a psychological account
using internal inference rules or proofs. That is different from using
logic as a standard of validity. A reasoner could implement a sound rule
poorly, use a model-search procedure soundly, or misinterpret a premise
before either procedure begins. Logical calculi do not by themselves
specify a natural-language parser, a resource allocation policy or which
of many valid conclusions will be reported. The source argues that model
construction handles meaning, background knowledge and counterexamples
more directly than a purely form-based psychological account. A rival
account may instead provide explicit semantic interpretation and control
procedures. Comparing the resulting systems requires their actual
predictions; declaring one of them ``logical'' is not a decisive test.

\begin{hbdefinition}{def:p19-r23-live-models}{Finite live possibilities and nonvacuous consequence}
Fix a finite set of Boolean atoms, let \(\Omega\) be its complete
valuations, and let \(\Gamma\) be interpreted propositional premises.
Write
\[
 \mathcal M(\Gamma)=\{v\in\Omega: v\models\gamma
                         \text{ for every }\gamma\in\Gamma\}.
\]
Ordinary classical consequence requires that every member of this set
satisfy the conclusion. A \emph{nonvacuous} consequence additionally
requires \(\mathcal M(\Gamma)\ne\varnothing\). A conclusion is a
\emph{live possibility relative to \(\Gamma\)} if at least one member
of \(\mathcal M(\Gamma)\) satisfies it. These are declared formal
benchmarks. An incomplete mental representation need not enumerate this
entire set.
For the two-atom examples, \(c\) denotes a circle's presence and \(t\)
a triangle's presence.
\end{hbdefinition}

\begin{hbtheorem}{thm:p19-r23-nonvacuity}{proved=annex:p19.r23.nonvacuity}{The satisfiable-premise boundary}
For satisfiable \(\Gamma\), ordinary and nonvacuous consequence agree.
For unsatisfiable \(\Gamma\), ordinary classical consequence holds for
every conclusion, whereas nonvacuous consequence holds for none and there
is no live possibility. In particular, \(\{c,\neg c\}\) is not a
basis for a nonvacuous prediction about \(t\).
\end{hbtheorem}

The source includes joint satisfiability in its proposed requirement for
right deduction. \Cref{thm:p19-r23-nonvacuity} states the effect explicitly;
it does not redefine classical consequence. Other proposed standards in
the chapter concern truth, survival or successful decisions. Logic,
probability, decision theory and game theory address different aspects of
these aims, and their psychological use is a further question. Strategic
reasoning among interacting agents is not developed in the selected
chapter. Likewise, accurate belief and efficient feasible action are not
the same criterion; compare \cref{sec:p17-rationality-8-5} and the
metareasoning distinctions at \cref{def:p19-r86-rationality-standards}.

\section{Deduction: alternatives, false clauses and counterexamples}
\label{sec:p19-r23-deduction}

Let \(c\) mean that a circle is present and \(t\) that a triangle is
present. The source's inclusive-disjunction display has three intuitive
models: a circle, a triangle, and both. Its fully explicit display gives
the corresponding truth assignments:
\[
 (c,\neg t),\qquad(\neg c,t),\qquad(c,t).
\]
The absence of a token in the intuitive display is not always an explicit
assertion of falsity. The distinction between leaving a clause unrepresented
and representing its negation is essential to the theory's predicted
illusory deductions.

\begin{hbtheorem}{thm:p19-r23-possibility}{proved=annex:p19.r23.possibility}{Live alternatives are not unrestricted modal consequences}
In the finite benchmark, each of \(c\), \(t\) and \(c\land t\)
is a live possibility relative to \(\{c\lor t\}\). Adding \(\neg c\)
removes the first and third possibilities. These judgments about the
premise-compatible set are not the unrestricted normal-modal inference
from \(c\lor t\) to \(\Diamond c\), nor the inference to
\(\Diamond(c\land t)\). Both modal inferences have countermodels even
on a one-world reflexive frame.
Under Boolean truth conditions, the assertion that exactly one of \(c,t\)
holds entails their inclusive disjunction, but the converse fails.
\end{hbtheorem}

The source proposes that a disjunctive assertion normally presents its
alternatives as genuinely open unless background information rules one
out. That interpretive default can support an ordinary possibility
response while an unrestricted modal entailment fails. It is not the
same as asserting that every conceivable modal logic invalidates the
inference. One must specify the interpretation and any added possibility
premises. The actual conditional-interpretation experiments at
\cref{sec:p19-conditional-interpretation-experiments} likewise separate
material truth conditions from assertions about possibilities.

The source's swan example assigns special importance to swans that are
white and excludes swans that are not white; nonswans are treated as
irrelevant to evaluating the generalization. Under classical interpretation,
``all swans are white'' and ``all nonwhite things are nonswans'' are
equivalent. Logic alone does not say how observing a nonwhite carnation
changes confidence in that hypothesis. Such a confirmation judgment
additionally requires sampling and background assumptions. A theory of
which examples people attend to is distinct from a theory of formal
equivalence or probabilistic evidence.
The chapter contrasts the confirmation puzzle with Einstein's
light-deflection prediction, whose specified observable consequence
makes a potential refutation intelligible. An empirical test still
evaluates auxiliary assumptions and measurements along with the
hypothesis, rather than logical equivalence alone.

The source also asks why a reasoner selects a novel, parsimonious conclusion
rather than an arbitrary logical consequence. Repeating that it rains,
or concluding that it rains and it rains, adds no information. From
``no teachers are poor'' and ``no parents are poor'', a response that no
relation between teachers and parents follows concerns a restricted
relational conclusion task: premises and tautologies still follow in
classical logic. Model theory postulates a conclusion-selection policy,
not the absence of those other consequences.

\begin{hbdefinition}{def:p19-r23-xor}{Exclusive alternatives in the source's illusion}
Here \(u\mathbin{\oplus}w\) means exactly one of the Boolean formulas
\(u,w\) is true. The source asks a reasoner to combine
\(c\mathbin{\oplus}t\) with \(c\mathbin{\oplus}\neg t\).
The first assertion's true alternatives can make a circle salient; the
second's true alternatives can do so too. Merely retaining that salient
token does not combine their full truth conditions.
\end{hbdefinition}

\begin{hbtheorem}{thm:p19-r23-xor}{proved=annex:p19.r23.xor}{A complete countercheck of the paired exclusive disjunctions}
The first assertion has precisely the valuations \((1,0),(0,1)\), and
the second has precisely \((1,1),(0,0)\), in coordinate order \((c,t)\).
Their conjunction is unsatisfiable. Accepting either \(c\) or
\(\neg c\) as a nonvacuous conclusion from the pair is therefore wrong
under these declared truth conditions.
\end{hbtheorem}

The theory predicts difficulty when false information must be recovered
and represented, and improvement when counterexamples are explicitly
searched for. \Cref{thm:p19-r23-xor} establishes the formal diagnosis; it
does not measure the prevalence of the predicted error. The 2021 chapter
describes its mSentential implementation as a model-generating procedure
with intuitive and deliberative modes. A program instantiation establishes
that a proposed procedure can be specified, not that it uniquely identifies
human cognition. Complete truth-table search is finite for a finite
propositional language, but its full valuation list has \(2^k\) entries
for \(k\) independent atoms. This size observation is not a proof that
all reasoning tasks require exhaustive enumeration or that every human
task is computationally intractable
\cite[pp.~148--150]{JohnsonLaird2021MentalModels}.

\subsection{An actual modal-inference experiment}

Hinterecker, Knauff and Johnson-Laird's Experiment~1 recruited 30
self-reported native English speakers through the UK Prolific Academic
service: 18 male and 12 female, mean age
28.4 years, SD 9.3. One participant who accepted every inference was
excluded, leaving 29. Each completed 12 inference trials, including six
inclusive-disjunction premises with a possible-conjunction conclusion and
four exclusive-disjunction premises with such a conclusion. Instructions
asked whether the conclusion followed necessarily if the premises were
true; ``possible'' was not further defined. Materials and item order were
counterbalanced/randomized as described in the study
\cite[early-online pp.~5--7]{HintereckerKnauffJohnsonLaird2016Modality}.

The reported acceptance percentages were 82\% for inclusive premises
and 10\% for exclusive ones; 26 of 29 participants accepted more of the
former than the latter, with one tie, reported binomial \(p<.00001\).
These results concern the administered wording and its interpretation.
The paper explicitly acknowledges that adding the premise that the two
alternatives can occur together changes the formal status of the
inclusive-to-possible-conjunction inference. Its logic comparison must not
be stripped of that qualification. The difference in acceptance does not
establish the unrestricted modal consequence refuted by
\cref{thm:p19-r23-possibility}, or distinguish every expanded mental-logic
account from every model account.

The two other inference types yielded 3\% acceptance of exclusive-to-
inclusive disjunction and 24\% of the converse. The first is truth-
preserving under Boolean semantics and the second is not, as in
\cref{thm:p19-r23-possibility}. The model theory predicts rejection of
both under its stronger interpretation that a disjunction asserts each
alternative's possibility: the inclusive conclusion introduces a joint
alternative excluded by the exclusive premise. This tests interpretation
as well as reasoning. Its predicted rejection must not be described as
invalidity of ordinary Boolean weakening.

Participants subsequently made six 0--100 probability judgments,
including four \emph{possibility statements}. Their mean values for those
four statements summed to 220\%. They were not mutually exclusive
cells of a binary joint distribution, so that sum alone is not a
normalization error. The paper's probability-validity comparison used
its stated assumption about joint possibility and analyzed 26 participants
with both relevant inference conditions; the corresponding difference
was marginal at reported \(p=.06\). This is not decisive rejection of
every probabilistic theory. The distinction between a probability of a
possibility statement and a probability assigned to one of four exclusive
outcomes matters both formally and empirically.

\section{Probability: possibility sets, weights and numerical estimates}
\label{sec:p19-r23-probability}

A model-based account of probability must add more than a list of possible
outcomes. The source proposes intuitive nonnumerical magnitudes of belief,
represented by an analog pointer or icon, and a deliberative mapping to
numbers and arithmetic. The pointer is a proposed computational
representation, not an observed anatomical device. Primitive averaging
of separately considered beliefs can violate joint-probability bounds;
numerical multiplication is not automatically correct either, because
independence must be justified.

The rival probabilistic account discussed in the chapter evaluates
inferences by their probability relations. In the one-premise monotonicity
benchmark, entailment entails that a conclusion is no less probable than
the premise. This simplified comparison must not be confused with a
complete multi-premise probabilistic consequence calculus. Further,
\(P(B\mid A)\), when defined, is not in general the probability of the
material conditional \(P(\neg A\lor B)\). The interpretation of a
natural-language conditional cannot be settled by interchanging them.

The source contrasts permission to enter a cinema now or in fifteen
minutes with a prediction that a storm could produce a flood or a tornado.
The permission example can invite separate permissions for each option;
the storm example does not assert that both effects must happen. Its
reasoner hears that bad weather makes flooding or a tornado possible,
learns that the weather is bad, and concludes that flooding is possible.
That response again depends on treating the named alternatives as open,
not merely on an unrestricted modal disjunction.
Permissions, epistemic possibilities and event probabilities have different
types. A number assigned to an event does not by itself encode a rule
about what an agent may do. The source's objection is that a probability
calculation alone leaves the construction of those possibilities and
intuitive estimates unexplained. A hybrid account could represent
possibilities as well as their weights; these examples do not prove
that no probabilistic account can include such semantic structure.

\begin{hbtheorem}{thm:p19-r23-joint}{proved=annex:p19.r23.joint}{Three assessments and a binary joint distribution}
For events \(A,B\), write \(a=P(A)\), \(b=P(B)\) and
\(q=P(A\cap B)\). Juxtaposition denotes intersection and a negated
event its complement. The four cells in the order
\((A B,A\neg B,\neg A B,\neg A\neg B)\) are
\[
 (q,\ a-q,\ b-q,\ 1-a-b+q).
\]
They always sum to one algebraically. For \(a,b\in[0,1]\), they
constitute a probability distribution exactly when
\[
 \max(0,a+b-1)\leq q\leq\min(a,b).
\]
For a union assessment \(u=P(A\cup B)\), substitute \(q=a+b-u\);
the admissible range is \(\max(a,b)\leq u\leq\min(1,a+b)\).
For a conditional assessment \(r=P(B\mid A)\) with \(a>0\),
substitute \(q=ar\) and retain the cell bounds. If \(a=0\), this
conditional is not determined by the ordinary ratio definition.
For an admissible joint distribution with \(a,b>0\), the converse
conditional \(r'=P(A\mid B)\) must equal \(ar/b\); if \(b=0\)
that converse ratio is undefined. When \(a\ne b\), their exact
arithmetic mean is neither an admissible intersection estimate nor
an admissible union estimate.
\end{hbtheorem}

Thus \(a=.3,b=.4,q=.5\) produces the cells
\((.5,-.2,-.1,.8)\), whose sum is one but which are not probabilities.
The source's dated 2024 Trump/Warren examples illustrate separate judgments
of events and their conjunction; they are not the actual unique-event
items used in the earlier study. An independently elicited fourth cell
can fail normalization, but the algebraically reconstructed four cells
from these three assessments fail through negativity, not through their
sum. The two assessment procedures have different possible failures.

\begin{hbtheorem}{thm:p19-r23-kernel}{proved=annex:p19.r23.kernel}{A structurally admissible triple and its null branches}
Choose \(a,r,s\in[0,1]\). The four numbers
\[
 (ar,\ a(1-r),\ (1-a)s,\ (1-a)(1-s))
\]
always form a probability distribution with \(P(A)=a\). If \(a>0\),
its \(P(B\mid A)=r\); if \(a<1\), its
\(P(B\mid\neg A)=s\). Every binary joint distribution has such a
representation. At \(a=0\) or \(a=1\), the unused conditional branch
can be chosen arbitrarily and is not identified by that joint distribution.
\end{hbtheorem}

The chapter uses this triple as a question about probabilistic expertise.
\Cref{thm:p19-r23-kernel} supplies the full answer and its boundary cases.
It is a reparameterization guaranteeing nonnegative cells, not a proof
that independently elicited human judgments are coherent or that a
respondent with this knowledge always uses it.

\subsection{The unique-events study and its limits}

Khemlani, Lotstein and Johnson-Laird's Experiment~1 tested 42 self-reported
native English speakers through Amazon Mechanical Turk, with a \$10
lottery as compensation. Each estimated a conjunction and its two
constituents for 16 unique-event materials in politics, sports and science.
Four judgment orders were administered within person, with four trials
per order and randomized presentation. The chosen event pairs had been
normed as increasing or decreasing one another's probability, not as
independent. Responses were entered on a 0--100\% slider; timing was
recorded without being mentioned in the instructions
\cite[early-view pp.~19--22]{KhemlaniLotsteinJohnsonLaird2015Naive}.

The study reported joint-distribution violations on 50\% of trials
overall, with 61\% when the conjunction was assessed first and 46\%
when last. The ordered trend was reported as Page's \(L=531\),
\(z=2.94,p<.003\). Twenty-eight of 42 participants gave conjunction
estimates between their two constituent estimates on at least half their
trials. This pattern is compatible with the proposed averaging process;
it is not direct measurement of an analog icon. Strategy classification
used a tolerance of 12.5 percentage points and an additional
between-constituents condition for averaging. Trial-order effects,
variable numerical judgments and the classification rule limit the
interpretation. A multiplication of marginals is admissible but represents
independence, which was not the normed relation for these items. Neither
the sample's observed proportion nor its reported \(p\)-values is a
population confidence interval or a universal frequency of irrationality.

Experiment~2 used a new sample of 41 from the same online population,
sixteen materials and two within-person orders: the inclusive union first
or last relative to its two constituent estimates. It reported
joint-distribution violations on 65\% of trials overall. Unlike the
conjunction study, order did not reliably alter violations (reported
\(p=.12\)) or additive-strategy classification (\(p=.20\)), contrary
to the theory's third prediction about arithmetic availability. Compound
estimation was nevertheless slower when first. Experiment~3 recruited
43 from the same population and used twelve materials, six with each
normed direction of influence. Its fixed order was \(P(A)\),
\(P(B\mid A)\), \(P(B)\), \(P(A\mid B)\). The first three
assessments violated cell admissibility on 24\% of trials; 42 of 43
participants produced at least one such trial. This fixed order did
not test the third prediction. The reported adjustment-strategy
comparison was post hoc, used a tolerance of 12.5 points, and allowed
an answer to meet both adjustment and Bayesian classifications
\cite[early-view pp.~23--28]{KhemlaniLotsteinJohnsonLaird2015Naive}.

The account's four study-level predictions concern shared judgments
across contents, patterned coherence violations, increased arithmetic
availability after constituent estimates, and faster repeated estimates
of information already considered. These are separable predictions,
not one indivisible confirmation. A valid converse-conditional
comparison additionally needs coherent inputs and nonnull conditioning
events, as in \cref{thm:p19-r23-joint}. None of these empirical results
is pooled here with the modal-possibility experiment.
The interpreted-event and language issues in
\cref{sec:p19-conjunction-language} remain possible competitors to an
unqualified cognitive-deficit diagnosis. A probability estimate also
becomes a betting commitment only under additional assumptions about
stakes, utility and two-sided prices; an answer on a slider is not
automatically such a commitment.

\begin{hbexample}{def:p19-r23-equiprobability}{Equal models require an extra weighting assumption}
The source's yellow-or-brown inclusive alternatives are yellow only,
brown only, and both. Assigning one equal chance to each gives probability
\(1/3\) to both. A different weighting of these same alternatives
need not give \(1/3\). Equal counting is a specified model of uncertainty,
not a consequence of the word ``or''.

In the source's disease example, let \(D\) be disease and \(S\) the
symptom, with chance counts
\[
 (DS,D\neg S,\neg DS,\neg D\neg S)=(3,1,2,4).
\]
Under the stipulated equal-chance sampling of these ten chances,
\(P(S)=5/10\) and \(P(D\mid S)=3/5\). This is a numerical
example, not prevalence data or a clinical diagnosis.
\end{hbexample}

\begin{hbtheorem}{thm:p19-r23-chances}{proved=annex:p19.r23.chances}{Chance counting and possible zero-mass outcomes}
The preceding chance calculation gives \(P(D\mid S)=3/5\).
For the source's exclusive assertion
\(R\mathbin{\oplus}(G\land B)\), the admissible Boolean triples
\((R,G,B)\) are exactly \(100,110,101,011\). The event
\(R\land B\) is possible because of \(101\), and has probability
\(1/4\) under a uniform measure on the four triples. It can have
probability zero under another probability measure on the same possibility
set. Consequently, possibility alone does not force a positive probability
without a support or weighting assumption.
\end{hbtheorem}

The red/green/blue-marble example illustrates the source's predicted
omission of a false clause: representing only red, or green-and-blue,
can conceal the red-and-blue case. A response of zero has an illusory
interpretation if the task grants positive weight to each represented
alternative; it is not an algebraic impossibility for every measure.
Similarly, a stipulated reference-class belief that incumbents often
win can alter an intuitive election estimate. Whether that belief
accurately describes a specified historical population is a separate
empirical question. A retrieved belief and verified reference-class
data are different inputs \cite[pp.~150--152]{JohnsonLaird2021MentalModels}.

\section{Induction, testing and defeasible interpretation}
\label{sec:p19-r23-testing}

The source's island example starts with ten observed overweight men
and compares generalizations that all inhabitants are obese or that
all are men. Background knowledge can make these extrapolations
differently plausible; neither follows deductively from the observations.
Similarity and contact can suggest causal hypotheses but
can also mislead. John Snow's cholera investigation and the Wright
brothers' attention to glider control serve as historical illustrations
of finding discriminating tests and causes, not controlled experiments
estimating a common mental-model effect
\cite[pp.~152--154]{JohnsonLaird2021MentalModels}.

For a fixed material conditional ``if it rains, the plants do not die'',
the additional premise that the plants died supports the modus-tollens
conclusion that it did not rain. For ``if it rains, it does not pour'',
combined with pouring and the background rule that pouring entails rain,
the premises cannot all be retained under that interpretation. A reasoner
may reject the first premise or reinterpret it as a defeasible
generalization. The latter response is premise revision, not a
counterexample to material modus tollens. Likewise, a failed bulb can
suggest an impurity or another failure mechanism without making the
suggested cause a deductive consequence. The suppression treatments in
\cref{sec:p19-conditional-reasoning-suppression} make the relevant
interpretation and information changes explicit.

Clinical illustrations require accurate causal content. Antibiotics do
not treat viral sore throats, and the symptom ``sore throat'' alone does
not establish a bacterial cause. Casual touching, hugging or sharing
dishes does not transmit HIV. These examples show why a plausible-looking
similarity or contact story must not be substituted for evidence about
the causal mechanism \cite{CDC2015Antibiotics,CDC2016HIVStigma}.

\begin{hbtheorem}{thm:p19-r23-cards}{proved=annex:p19.r23.cards}{Complete testing of the four displayed cards}
Suppose each of four cards has one letter and one number. Their exposed
faces are E, F, 2 and 3. The source's rule says that any card with E has
2 on its other side. With no other information about hidden faces,
turning E and 3
is sufficient and individually necessary to decide conclusively whether
these four cards satisfy the rule. Neither F nor 2 can be a counterexample
regardless of its other face. This conclusion uses a specified complete
verification goal, not an unrestricted theorem about optimal information
search under every prior and cost.
\end{hbtheorem}

The source contrasts the familiar tendency to seek a matching positive
case with explicit search for a falsifying case. Whether a card could
refute a material rule, whether inspecting it increases a Bayesian
information score, and whether inspecting it is a socially relevant
check are different objectives. The model account predicts dependence
between selections, for example choosing 2 together with E as a possible
example, rather than treating each card choice as an independent response.
It also predicts that task interpretation and content affect the
availability of potential counterexamples. The chapter's discussion of a
large selection-task literature connects those predictions to content
and familiarity effects. Such effects do not license transferring an error
frequency from an abstract-card sample to every hypothesis-testing situation.

Ragni, Kola and Johnson-Laird's 2018 review assembled 228 experiments
reporting the four canonical selections: 104 abstract, 44 everyday and
80 deontic. This is a literature database, not a newly recruited common
sample. Its dependence analysis used the smaller set of 99 experiments
reporting all sixteen possible card subsets: 43 abstract, 30 everyday
and 26 deontic. For each experiment, 10,000 independent-card simulations
preserved its sample size and the observed marginal probabilities of
selecting each card. Mean observed versus simulated selection entropy,
in bits, was 2.33 versus 2.60 for abstract, 2.43 versus 2.75 for everyday,
and 1.57 versus 1.80 for deontic tasks. All three paired Wilcoxon
comparisons reported \(p<.001\). This tests independent selection
against redundancy in complete response subsets; dependence is not
established by that analysis for all 228 experiments, nor does it
uniquely identify mental models as its cause
\cite[pp.~783--786, Tables~1--3]{RagniKolaJohnsonLaird2018Selection}.

Within the canonical selections, counterexample choices were 19\%,
28\% and 64\% for the abstract, everyday and deontic collections.
The content and abstract-instruction rank comparisons reported
chi-square statistics of 54.6 and 62.3, respectively, each with two
degrees of freedom and \(p<.001\).
Those percentages concern the four canonical response categories,
not the prevalence among every response in a target population.
The review also retains studies in which content or instruction
changes had no reliable effect. Its model-fitting comparison used
three model-theory parameters against a rival simplified from ten
to four parameters to fit canonical selections. The authors acknowledge
that the pruning can be disputed and that better fit in this one task
does not establish superiority of an entire theory. Unexplained insight
parameters remain a mechanism gap. Neither the fit nor the content
contrast supplies a randomized training effect or warrants rejecting
every probabilistic account under every interpretation
\cite[pp.~785--786,789--792, Tables~4,6]{RagniKolaJohnsonLaird2018Selection}.

\subsection{Repeated selection of concrete shapes}

The 2021 chapter illustrates the repeated task using a triangles-are-white
generalization. The original 1970 experiment below used triangles-are-black.
Swapping colors leaves the verification structure unchanged if the
candidate-counterexample roles are swapped as well, but the two
descriptions are not identical administered materials.

Johnson-Laird and Wason's 1970 Experiment~I tested 28 first-year psychology
undergraduates at University College London, allocated in rotation to
four groups of seven. The groups crossed conditional versus quantified
wording of a triangles-are-black rule with instructions to prove it true
versus false. Participants requested unreplaced objects, one at a time,
from a white or a black box, each known to contain 15 objects; only circles
and triangles were used. The black box contained ten triangles and five
circles. White boxes contained fifteen circles in the prove-true
conditions, and fourteen circles followed by one triangle in the
prove-false conditions. Predetermined black-box ordering initially gave
five triangles. Incorrect premature proof claims were rejected and
participants told to continue \cite[pp.~50--53]{JohnsonLairdWason1970Insight}.

Every participant eventually exhausted the white box and no participant
exhausted the black box. This supports eventual task-specific recognition
that a white triangle is the possible falsifier. The mean number of
unnecessary black objects requested was 4.29 in the prove-true condition
and .86 in the prove-false condition, with reported one-tailed
\(p=.0002\). The main conditional-versus-quantified difference was not
significant; a secondary timing-of-first-black-request measure, after
excluding initial perfect responders, did differ. Eventual correct
selection with feedback is not immediate insight without assistance.

The same paper's Experiment~II tested 24 native-English-speaking
psychology undergraduates unfamiliar with formal logic, using eleven
four-dot diagrams and six rules. The task was to decide whether a rule
described the diagrams in an envelope by asking about diagram locations.
Five participants consistently selected positive and negative cases and
four showed complete insight from the start; their data were discarded
for the paper's focal analysis, leaving 15. Some of those participants
lost the task's negative-instance focus for more complex rules,
especially a disjunctive rule. This within-paper contrast qualifies the
claim of permanent insight from repetition; it is not a clean causal
estimate of working-memory capacity or a measure of all reasoning
competence \cite[pp.~54--60]{JohnsonLairdWason1970Insight}.

\section{Abduction: causal explanations and algorithms}
\label{sec:p19-r23-abduction}

The source distinguishes restoring consistency by changing a belief from
constructing an explanation of why an apparent conflict arose. Friends
who leave a cafe to collect a car, promise to return in ten minutes and
remain absent after twenty invite hypotheses such as trouble starting
the car, an accident or a different route. None is entailed by lateness.
A causal account adds a mechanism and can support new tests; merely
deleting the expected-return belief need not do so. Conversely, an
invented explanation can hide an inconsistency rather than resolve it
\cite[pp.~154--155]{JohnsonLaird2021MentalModels}.
In the cafe episode, waiting avoided a further risk: searching could
separate the two groups if the friends returned meanwhile. The source
reports that the car eventually arrived after needing a tow to start.
That outcome is one historical confirmation of the proposed explanation,
not an experiment estimating the reliability of such waiting decisions.

The linked 2004 model account separates detecting a conflict, choosing
which belief to revise, and explaining its origin. It first tries to
construct a compatible model. When a new fact conflicts with the model,
it checks a particularly dubious belief or an explicit conflict before
using a mismatch between the representation and the fact to guide
retraction. Causal knowledge can then extend the facts into a candidate
explanation, whose consequences can alter other beliefs. The order is
not rigid: constructing an explanation may determine which belief is
revised. No arbitrary deletion criterion alone performs all three
operations \cite[pp.~640--642]{JohnsonLairdGirottoLegrenzi2004Consistency}.

That paper also reports probability rankings of causal explanations
for twenty scenarios. Cause-and-effect chains ranked above their
separate constituents and above a noncausal conjunction control; an
initial common scenario order was followed by a four-order replication.
Its report supplies neither the samples nor effect uncertainty for
these rankings. They must not become a population rate. With unchanged
event meanings and one probability law, a conjunction cannot exceed
either constituent (\cref{thm:p19-conjunction-benchmark}); a different
interpretation of the explanation would need separate justification
\cite[pp.~654--655]{JohnsonLairdGirottoLegrenzi2004Consistency}.

Khemlani and Johnson-Laird's 2011 pilot used 29 paid Mechanical Turk
participants without logic training, each freely answering what followed
from twelve conditional-plus-counterexample scenarios. Explanations were
coded as adding new causal information; two independent raters agreed on
80\% of trials and reconciled differences. Added-information responses
occurred on 69\% of trials, direct refutations on 31\%, with reported
Wilcoxon \(z=3.15,p<.005,\delta=.67\). Three further online samples
of 21, 17 and 25 participants respectively chose the most probable
supplied conclusion, ranked four conclusions, or rated one per trial.
Each design favored explanations qualifying a generalization, but the
ordering of other responses differed. In its proposed implementation,
nonfiring after a trigger pull retrieves possible disabling conditions,
then possible causes of them. The study does not verify those explanations
as facts, establish that they preceded every belief revision, or rule out
probabilistic accounts allowing defeasible interpretation. Causal content,
conditional wording and the supplied response menus restrict the
comparison \cite[early-online pp.~3--9]{KhemlaniJohnsonLaird2011Explain}.

Their 2012 paper tested a different prediction: explaining first can
reduce later detection of a strictly interpreted inconsistency. Its five
studies recruited Mechanical Turk samples of 40, 53, 40, 38 and 25.
The first specified North American subscribers, self-reported native
English speakers without logic expertise; subsequent samples used the
same platform. Experiment~1 assigned 20 participants per task order;
inconsistent-pair detection was 64\% with evaluation first and 27\%
with explanation first. Later designs added ``always'', explicitly
requested literal evaluation, compared surprise judgments or their
justifications, and finally made the conditions within-person. The last
study reported inconsistent-pair detection of 60\% with no prior task,
56\% after surprise, and 40\% after explanation, with its interaction
reported as \(z=2.08,p<.025,\delta=.23\). The paper acknowledges
unresolved explanation-content and mechanism questions. Its interpretation
is a change toward exception-tolerating premises, not new satisfaction
of unchanged classical formulas. These studies are distinct from
Lombrozo's supplied-explanation preference tasks at
\cref{sec:p19-explanatory-preferences}
\cite[pp.~486--490]{KhemlaniJohnsonLaird2012HiddenConflicts}.

A cobra's venomous bite can cause death, but it does not guarantee
death in every case.
Survival may prompt investigation of effective treatment, injected venom
and the original diagnosis. A tourniquet or suction is not an established
rescue mechanism: CDC/NIOSH explicitly advises against both. A program's
stored causal story is not medical evidence
\cite{NIOSH2019VenomousSnakes}. More generally, failing to find a model
of a premise set does not establish that the set is inconsistent unless
the search is complete. In the source's mSentential illustration, the
program chooses a stored candidate at random, then searches for its cause
to extend the causal chain. That selection differs from a person's
assessment of candidate likelihoods. It also constructs a counterfactual
description of what its assumed mechanism would produce without the
selected cause. The counterfactual depends on the causal model, not
co-occurrence alone. The procedure's existence is not a comparative benchmark
showing that humans outperform every computer explanation system.

\begin{hbdefinition}{def:p19-r23-railway}{Word semantics for the railway example}
A state is \(L[S]R\), where \(L,S,R\) are finite words of car
labels (which may repeat) and \(\epsilon\) is the empty word. Words are written in the
physical left-to-right order used by the source's diagrams. The allowed
movements are: transfer a nonempty terminal block \(T\) of \(L\) to
the beginning of \(S\), transfer a nonempty initial block of \(S\)
to the end of \(L\), or transfer a nonempty terminal block of \(L\)
to the beginning of \(R\). Each transfer preserves the internal order
of its block. There is no direct siding-to-right move, no right-to-left
move and no leap over another car. This specifies the task's moves,
not a new logical calculus.
\end{hbdefinition}

\begin{hbmethod}{def:p19-r23-reversal}{The chapter's reversal procedure}
For a nonempty initial word \(a_1\cdots a_n[\epsilon]\epsilon\),
transfer all but \(a_1\) to the siding if \(n>1\). Transfer
\(a_1\) to the right. While the siding is nonempty, transfer its
first car to the left and then that car to the right. For the empty input,
do nothing. The first instruction is a bulk move when required; the loop
uses two one-car transfers per iteration.
\end{hbmethod}

\begin{hbtheorem}{thm:p19-r23-reversal}{proved=annex:p19.r23.reversal}{A full correctness invariant for every finite train}
The method terminates and produces
\(\epsilon[\epsilon]a_n\cdots a_1\) for every finite input word,
including the empty word. For \(n\geq1\), after \(k\) cars have
been placed on the right, the loop-boundary state is
\[
 \epsilon[a_{k+1}\cdots a_n]a_k\cdots a_1,
       \qquad 1\leq k\leq n.
\]
The procedure uses one move when \(n=1\), and \(2n\) transfers when
\(n>1\), counting the initial bulk transfer as one. This is a
correctness and stated-move-count result, not a proof of globally shortest
program length.
\end{hbtheorem}

For ABCD, the source's trace is
\[
 ABCD[\epsilon]\epsilon,\quad A[BCD]\epsilon,\quad
 \epsilon[BCD]A,\quad B[CD]A,\quad\epsilon[CD]BA,
 \quad\ldots,\quad\epsilon[\epsilon]DCBA.
\]
The program mAbducer derives informal procedures by simulating instances,
finding repeated move sequences and translating them into loop
descriptions. A for-loop can encode a fixed number of repetitions as a
function of train length; a while-loop can use the observed empty-siding
condition. Discovering a general loop from a few simulations is not by
itself proof of correctness for every length. The invariant in
\cref{thm:p19-r23-reversal} supplies that additional step.

The source distinguishes performing a rearrangement, inventing a procedure
for it and predicting the result of a supplied procedure. It predicts
that the first task's difficulty depends on moves and cars per move,
whereas algorithm formulation depends on the complexity of describing
the procedure. Reversal, palindromic regrouping, parity sorting and the
inverse parity rearrangement provide different structural demands.
Performing many moves can therefore be compatible with describing a
short loop. The actual mental-rotation benchmarks at
\cref{sec:p19-mental-rotation} illustrate another distinction between
behavioral timing and identifying a unique internal simulation process.

Khemlani and colleagues' 2013 algorithm-abduction experiment tested
20 nonprogrammers from the Princeton University undergraduate population.
After five practice problems with the railway interface,
the interface was switched off. Participants wrote informal procedures
for reversal, palindrome and parity-sort problems for eight-car trains and
for trains of arbitrary length, in two counterbalanced blocks. The study
reported 92\% correct for the fixed-length problems and 52\% for the
arbitrary-length problems, with Wilcoxon \(z=3.29,p<.001\).
Loop classifications for arbitrary-length solutions used the 18
participants who produced at least one correct such solution. Twelve
used more while-loops than for-loops, with three ties and reported
binomial \(p<.02\). The difference is consistent with the proposed
simulation account but does not establish that for-loops require solving
equations in every human task or that observed descriptions uniquely
identify their cognitive origin
\cite[main paper pp.~3--6 and SI Text S3--S4]{KhemlaniEtAl2013Kinematic}.

In Bucciarelli and colleagues' 2016 Experiment~3, 34 children from the
same Turin primary-school population as the preceding studies were
recruited; their mean age was 10 years 5 months, range 9 years 10 months
to 11 years 5 months, with 13 male and 21 female. Two unclear-description
cases were excluded, leaving 32. Each described four six-car
rearrangement procedures twice without moving the cars, once able to
gesture and once with arms crossed inside a cloth muff. Block order and
letter-versus-number car labels were counterbalanced, with randomized
order within blocks. Coding used speech and gesture, independently
judged and reconciled before analysis
\cite[early-online pp.~12--17]{BucciarelliEtAl2016Algorithms}.

Correctness was reported as 65\% with gestures and 52\% without,
a 13-\emph{percentage-point} difference, with Wilcoxon
\(z=2.37,p<.01\) and Cliff's \(\delta=.23\). Mean times for
correct descriptions were reported as 40 and 44 seconds respectively.
The paper discusses carryover, the conspicuous muff manipulation,
gestures' communicative role and unobserved movements as limitations.
This experiment supports a task-specific gesture contribution; neither
the difference nor spontaneous gestures alone proves the mental-model
mechanism or establishes a general educational prescription.

\begin{proofboundary}{def:p19-r23-program-length}{Description length is relative to a language}
A count of instructions or of bits in a chosen implementation is a
declared description-length measure. It need not be the length of the
shortest program producing the same function. The studies' measured Lisp
description lengths are therefore upper bounds in that representation,
not computed exact Kolmogorov complexities. General uncomputability of
shortest descriptions does not prevent a specific lower-bound or
minimality proof under explicit finite primitives. Likewise, a
while-loop with an unbounded stopping condition and a for-loop with a
fixed primitive bound belong to different restricted languages; ordinary
programming-language loop names alone do not determine computational power.
\end{proofboundary}

\section{Using model construction and evaluating its limits}
\label{sec:p19-r23-use}

The source's practical model method, associated with Victoria Bell,
externalizes alternatives in columns, represents their true and false
components explicitly, combines them with new information and removes
incompatible cases. A reasoner can then ask whether a proposed conclusion
holds throughout the remaining possibilities or whether a counterexample
remains. For finite interpreted propositional tasks the complete
procedure is justified by \cref{def:p19-r23-live-models}; its correctness
depends on an exhaustive and accurate representation. Training people
to use it is a different empirical proposition, not a consequence of
that formal justification.

The chapter's reported training benefit can be traced to Johnson-Laird's
2006 book, which describes evaluation of given conclusions in two blocks
of reasoning problems and three groups: unaided in both blocks, taught
the model method after the first block, and taught before the first block.
The book reports about two-thirds correct and mean evaluation time of
24 seconds without instruction, versus 95\% and 15 seconds with it.
This design is not a single individually paired 66\%-to-95\% pre/post
comparison. The accessible passage supplies no participant count,
recruitment, allocation procedure, item count, inferential test,
confidence interval or delayed-transfer measure. Its citation leads to
Bell's 1999 unpublished dissertation, whose actual methods and results
have not been verified here. These numbers therefore remain a
book-reported result with an unresolved underlying-study evidence
boundary \cite[pp.~288,290]{JohnsonLaird2006HowWeReason}.

The same book separately describes an easier experiment without pencil
and paper, using problems with biconditional or exclusive-disjunction
premises. It reports 80\% correct unaided versus 93\% after instruction.
Those are not the first experiment's accuracy or timing results; no
missing sample size or uncertainty estimate can be inferred from the
percentages \cite[p.~290]{JohnsonLaird2006HowWeReason}.

The comparison with conventional logic training concerns a particular
selection-task study, not the whole effectiveness of logic education. Cheng and
colleagues' Experiment~2 analyzed 53 students completing both tests in
two University of Michigan introductory logic classes, at Ann Arbor
and Dearborn. About forty lecture hours covered propositional inference,
truth tables and proofs without selection-task instruction. Four matched
selection problems were given in the first and final weeks, with
counterbalanced test sets and no intervening feedback. The percentage
of conditional problems with at least one error fell by three points,
with a reported 95\% comparison-interval half-width of seven points;
this did not establish improvement or equivalence. Erroneous selection
of the consequent did decrease. There was no randomized untreated-class
comparison. Their laboratory study also found benefits from combining
rule instruction with application examples. Pragmatic-schema training
is therefore a relevant competitor, not evidence that explicit rules
can never help \cite[pp.~306--309]{ChengEtAl1986Training}.

Neither this classroom comparison nor a model-method improvement
establishes transfer to every reasoning domain. The chapter's survey
remark about psychologists and
the admissible-triple question supplies no sample, sampling frame or
uncertainty estimate; it cannot support a population prevalence. The
formal answer is nevertheless fully provided at
\cref{thm:p19-r23-kernel}.

As a working method, state the interpretation, list the relevant
alternatives, make omissions and negations explicit, and identify the
kind of conclusion sought. For a necessity claim, attempt a
counterexample; for a possibility claim, exhibit a compatible case; for
a probability claim, declare weights and check joint admissibility; for
an explanation, distinguish a candidate mechanism from a verified one;
for an algorithm, establish an invariant and a stopping condition.
Resource costs determine whether full enumeration is feasible, while
confidence and decisions to stop are additional processes, treated at
\cref{sec:p19-rationality-8-6}. This method is not permission to treat
one imagined case as exhaustive evidence.

The strongest comparison between model, mental-logic and probabilistic
accounts uses common interpreted tasks, independently specified
predictions and evidence that discriminates their mechanisms. Each can
be expanded with semantic and resource assumptions, sometimes reducing
the original contrast. The model account's false-information omission,
explicit-counterexample search, numerical-order effects and
rearrangement-versus-algorithm distinction are substantive predictions;
they are not established merely by a reader finding the examples
intuitive. The source's proposed rationality criterion joins accurate
representation, truth, parsimony and explanatory power. These aims still
require choices about evidence, acceptable risk and the cost of further
search. Empirical frequency cannot settle those normative choices.
